\documentclass[12pt]{article}

% Shared preamble of the paper and its internet appendix.
\usepackage[margin=1in]{geometry}
\usepackage{amsmath,amssymb}
\usepackage{booktabs}
\usepackage{colortbl}
\usepackage{array}
\usepackage{caption}
\usepackage{float}
\usepackage{tikz}
\usepackage{graphicx}
\usepackage{setspace}
\usepackage{xcolor}
\usepackage{enumitem}
\usepackage{xr-hyper}
\usepackage{hyperref}
\usepackage[backend=biber,style=authoryear,natbib=true,maxcitenames=4,maxbibnames=20,
            uniquename=false,uniquelist=false,giveninits=false,dashed=true]{biblatex}
\setcounter{biburllcpenalty}{7000}
\setcounter{biburlucpenalty}{8000}
\setcounter{biburlnumpenalty}{7000}
\addbibresource{references.bib}
\addbibresource{references-erc.bib}

% Journal of Finance-style reference list: "Author, First, Year, Title, Journal Volume, Pages."
\DeclareNameAlias{sortname}{family-given/given-family}
\renewcommand*{\revsdnamepunct}{\addcomma}
\DeclareDelimFormat{finalnamedelim}{\ifnumgreater{\value{liststop}}{2}{\finalandcomma}{}\addspace\bibstring{and}\space}
\renewcommand*{\labelnamepunct}{\addcomma\space}
\renewbibmacro*{date+extradate}{%
  \iffieldundef{labelyear}{}{\printtext{\printlabeldateextra}}}
\DeclareFieldFormat*{title}{#1}
\DeclareFieldFormat*{journaltitle}{#1}
\DeclareFieldFormat[article]{volume}{#1}
\DeclareFieldFormat{pages}{#1}
\DeclareFieldFormat{url}{\url{#1}}
\renewbibmacro*{in:}{}
\renewbibmacro*{volume+number+eid}{\printfield{volume}}
\renewcommand*{\bibpagespunct}{\addcomma\space}
\AtEveryBibitem{\clearfield{number}\clearfield{doi}\clearfield{issn}\clearfield{isbn}%
  \clearfield{urlyear}\clearfield{urlmonth}\clearfield{urlday}\clearfield{eprint}%
  \ifentrytype{article}{\clearfield{url}}{}%
  \ifentrytype{misc}{\clearfield{note}\clearfield{addendum}}{}%
  \ifentrytype{online}{\clearfield{note}\clearfield{addendum}}{}}
\setlength{\bibhang}{1.5em}
\renewcommand*{\bibfont}{\small}
\DeclareDelimFormat[bib]{finalnamedelim}{\addcomma\space and\space}
\DeclareDelimFormat[bib]{nameyeardelim}{\addcomma\space}
\AtBeginBibliography{\renewcommand*{\newunitpunct}{\addcomma\space}}

\newcommand{\gitrepo}{https://github.com/GF-boop/beyond-stocks}
\newcommand{\gitlink}[2]{\href{\gitrepo/blob/main/#1}{#2}}
\newcommand{\gitdirlink}[2]{\href{\gitrepo/tree/main/#1}{#2}}

\hypersetup{pdftitle={Beyond All-Equity: Levered Diversification over the Lifecycle},pdfauthor={Guillaume Flament},colorlinks=true,citecolor=blue!55!black,linkcolor=blue!55!black,urlcolor=blue!55!black}
\renewcommand{\thetable}{\Roman{table}}
\captionsetup[table]{labelsep=colon,font=normalsize,skip=4pt}
\captionsetup[figure]{labelsep=period,font=small,justification=justified,labelfont=bf}
\setlength{\emergencystretch}{3em}
\setlength{\parindent}{1.5em}
\setlength{\parskip}{0pt}
\onehalfspacing

\externaldocument{internet-appendix}

% Exhibits follow the text in the working-paper version. Set \submissiontrue to
% move every table and figure after the references, in order of appearance.
\newif\ifsubmission
\submissionfalse
\makeatletter
\newcommand{\exhibitlist}{}
\newcommand{\exhibit}[1]{\ifsubmission\g@addto@macro\exhibitlist{\input{#1}}\else\input{#1}\fi}
\makeatother

\begin{document}

\begin{titlepage}
\begin{center}
{\large\bfseries BEYOND ALL-EQUITY: LEVERED DIVERSIFICATION\\[2pt] OVER THE LIFECYCLE}\\[12pt]
Guillaume Flament\footnote{Cayas. Email: \href{mailto:guillaume.flament@cayas.fr}{guillaume.flament@cayas.fr}.}\\[8pt]
September 22, 2026\footnote{The replication code, data, and numerical results are available at
\url{\gitrepo}. The internet appendix is available as a separate document. The author used OpenAI Codex
and Anthropic Claude Code for code development, reproducibility checks, and manuscript editing, and is
responsible for the research design, interpretations, code, and final manuscript.}
\end{center}
\vspace{6pt}
\begin{center}\textbf{Abstract}\end{center}
\begin{quote}\small
We study whether the case for all-equity lifecycle portfolios survives when households can borrow and
invest in global bonds, gold, and managed futures. Without leverage, a multi-asset strategy lowers the
probability of financial ruin but not expected utility: it requires a 12.35\% savings rate to match an
all-equity strategy saving 10\%. Levering equities alone adds little utility and raises ruin. Levered to
the volatility of unlevered equities, the multi-asset strategy matches the all-equity utility with a
6.35\% savings rate, lowers the probability of ruin from 7.0\% to 2.3\%, and, at the same savings rate,
fails to fund the all-equity retirement income less often (3.2\% versus 7.0\%). The gains come mainly from
managed futures, which pay off in equity crises. A strategy of stocks and managed futures alone does
better still on every measure, whereas bonds lower utility when they replace stocks.

\medskip
\noindent JEL classifications: C15, D14, G11, G17, G51\\
Key words: Lifecycle asset allocation, leverage, managed futures, retirement savings, risk parity
\end{quote}
\begin{quote}\small
\textbf{Key findings}
\begin{itemize}[leftmargin=1.2em,itemsep=1pt,topsep=2pt]
\item Levered to the volatility of equities, a multi-asset portfolio gives a couple the expected utility
of an all-equity portfolio with 6.4\% of income saved instead of 10\%, and cuts its risk of running out of
money from 7.0\% to 2.3\%.
\item Managed futures drive the gains; stocks plus managed futures do best. Gold helps since 1970, when
its price is free, but not when it was administered. Bonds lower utility when they replace stocks.
\item Preserving capital and funding a given income differ: risk parity at 175\% runs out of money in 3.8\%
of lifetimes but fails to pay the all-equity income in 14.2\%.
\end{itemize}
\end{quote}
\end{titlepage}

\section{Introduction}

\citet{AnarkulovaCederburgODoherty2025} challenge two central tenets of lifecycle investing. In a
lifecycle model with labor income risk, Social Security income, and longevity risk, a US couple that
chooses among domestic stocks, international stocks, bonds, and bills optimally holds about one-third
domestic stocks and two-thirds international stocks at all ages, with virtually no fixed income. This
all-equity strategy dominates balanced portfolios and target-date funds in building wealth, preserving
capital, and generating bequests. The conclusion survives leverage. When their couple can borrow at a
modest spread above the bill rate, it levers the all-equity portfolio rather than a balanced portfolio
of the kind advocated by \citet{Asness1996}. \citet{Asness2024} objects that diversified portfolios
should be compared with equities at equal risk, after levering them, not at equal capital. We test
this objection with a broader set of diversifiers.

We add two asset classes that have historically paid off during equity crises, gold and managed futures,
alongside currency-hedged global government bonds. These assets have lower expected returns than stocks,
so an unlevered household that holds them gives up return. Leverage is therefore the natural companion to
diversification: a household that can borrow can hold a more diversified portfolio and scale it up to an
equity-like level of risk. This is the portfolio logic behind the leverage-aversion argument of
\citet{Asness1996}, \citet{FrazziniPedersen2014}, and \citet{AsnessFrazziniPedersen2012}.

We keep the household problem of \citet{AnarkulovaCederburgODoherty2025} and implement it at an annual
frequency on a public dataset of 16 developed countries from 1927 to 2025. A US couple saves 10\% of its
labor income from age 25, withdraws 4\% of its retirement wealth in real terms from age 65, receives
Social Security benefits, and has utility over retirement consumption and bequest. We compare an
all-equity strategy, a proportional strategy that holds stocks, global bonds, gold, and managed futures in
fixed proportions of 40\%, 27\%, 17\%, and 17\%, and a risk parity strategy, financed at the local bill
rate plus 0.30\%. Following the objection of \citet{Asness2024}, our main comparison levers each
diversified strategy to the volatility of the unlevered all-equity strategy; we also report exposures
from 100\% to 200\% of wealth.

We report three measures, which answer different questions (Table~\ref{tab:measures}). The equivalent
savings rate is the savings rate that gives a strategy the expected utility of the all-equity strategy
saving 10\%. The probability of financial ruin applies the 4\% rule to the strategy's own wealth and does
not depend on the savings rate. The path-matched income failure is the probability of not paying the
income that the all-equity strategy would have funded on the same simulated path.

Figure~\ref{fig:ladders} summarizes our main results. Without leverage, the proportional strategy lowers
the probability of ruin from 7.0\% to 5.0\% but does not improve expected utility: a couple using it must
save 12.35\% of its income to match an all-equity couple saving 10\%. Levering the all-equity strategy adds
little utility and raises its ruin probability at every step. Levered to the volatility of unlevered
equities, at 166\% exposure, the proportional strategy improves on all three measures. It matches the
all-equity utility with a 6.35\% savings rate, has a ruin probability of 2.3\%, and, saving 10\%, fails to
fund the path-matched income in 3.2\% of lifetimes, compared with 7.0\%.
\exhibit{figures/erc/fig_ladders.tex}

These three numbers should not be combined. The equivalent savings rate is a utility measure: a couple
that saves only 6.35\% in the proportional strategy obtains the same expected utility, not the same
income, and it fails to fund the path-matched income in 13.1\% of lifetimes, most often in paths where
equities did well and that income is high. A strategy of stocks and managed futures alone, levered to
132\%, does better on every measure: it matches the all-equity utility with a 5.22\% savings rate and, even
at that savings rate, funds the path-matched income more reliably than the all-equity strategy.

These findings are consistent with the portfolio logic of leverage aversion but not with its usual
source. In the argument of \citet{AsnessFrazziniPedersen2012}, low-risk assets, bonds in particular,
offer higher risk-adjusted returns, and an unconstrained investor should lever them. In our data, global
bonds have a Sharpe ratio of 0.29, no higher than the 0.31 of stocks (Table~\ref{tab:sleeves}). The gains
come instead from managed futures, which have a Sharpe ratio of 0.49 and gain 8.4\% on average in the
country-years in which stocks lose most. Removing managed futures from either diversified strategy raises
its equivalent savings rate, whereas removing bonds lowers it. Gold contributes after 1970, when its
price was free.

The choice of retirement outcome also matters in its own right. Under the 4\% rule, a strategy with low
expected returns can preserve its capital simply because it promises a small income. At 175\% exposure,
the risk parity strategy has a ruin probability of only 3.8\%, but it fails to fund the path-matched
income in 14.2\% of lifetimes, twice the rate of the all-equity strategy. Levered to equity volatility,
at 265\% exposure, risk parity does better than the proportional strategy, but this advantage is
fragile.

Our conclusions rest on two conditions: cheap financing and the returns of our managed-futures series,
which we build from public futures data back to 1927. Both are within reach of individual investors today.
At equal volatility, the proportional strategy keeps its saving advantage up to a borrowing spread of about
2.75\% and its income advantage up to about 1.3\%. Individual investors do not need a margin loan to borrow at such rates. Box spreads on S\&P 500 index
options, which trade in ordinary brokerage accounts, have been priced at 0.2\% to 0.5\% above comparable
SOFR or Treasury rates \citep{CMEBoxFinancing2024,WangBoxFinancing2025},
and the financing embedded in S\&P 500 futures, which return-stacked funds use, cost 0.3\% on average over
2021 to 2023 \citep{DEShaw2025}. Our series earned 0.7 to 2.0 percentage points per year more than
commercial trend-following indexes after 2000, yet even with a penalty of 2\% per year on its returns, the
proportional and stocks-and-managed-futures strategies still improve on the all-equity strategy in
utility, ruin, and income funding. Funds that replicate these indexes are available to individual
investors, including as UCITS exchange-traded funds \citep{iMGPDBMF2025}. The results also hold across
resampled calendar histories, managed-futures signals, samples, preferences, and savings policies.

We contribute to the normative literature on lifecycle investment strategies
[e.g., \citet{CoccoGomesMaenhout2005}; \citet{Gomes2020}; and \citet{AnarkulovaCederburgODoherty2025}].
Our primary contribution is to separate the value of diversification from the value of leverage in a
lifecycle model with labor income, Social Security, and longevity risk. Most lifecycle studies restrict
households to stocks, bonds, and bills and rule out borrowing.\footnote{\citet{DavisKublerWillen2006}
show that realistic borrowing costs reduce the demand for equity over the life cycle.
\citet{AyresNalebuff2010,AyresNalebuff2013} propose leveraged equity positions for young investors to
diversify across time. At the lowest spread of \citet{AnarkulovaCederburgODoherty2025}, their couple
borrows 100\% of wealth and holds 15\% in bonds.} We show that the gains from leverage depend on the assets
being levered, and that the gains from diversification depend on leverage. Our second contribution is to
distinguish capital preservation from the ability to fund a given retirement income. Finally, we use a
panel built entirely from public data, and we release the data and code.

\section{Leverage and diversification in lifecycle investing}

In the models of \citet{Merton1969} and \citet{Samuelson1969}, a sufficiently risk-tolerant investor
borrows to hold more of the risky asset. Later lifecycle models typically rule out borrowing or make it
costly. \citet{CoccoGomesMaenhout2005} impose borrowing constraints, and \citet{DavisKublerWillen2006}
show that the spread between borrowing and lending rates lowers optimal equity holdings.

A separate literature argues that borrowing constraints distort the choice of risky assets, not only
the level of risk. \citet{Asness1996} shows that a levered stock-bond portfolio historically delivered
equity-like returns with lower volatility than an all-equity portfolio. \citet{FrazziniPedersen2014}
and \citet{AsnessFrazziniPedersen2012} provide a theory and evidence: investors who cannot or will not
borrow overweight risky assets, which lowers the risk-adjusted returns of those assets. Risk parity
portfolios, which equalize risk contributions across asset classes
\citep{MaillardRoncalliTeiletche2010}, are the standard institutional application of this logic. Their
historical performance depends on the financing cost and the sample period
\citep{AndersonBianchiGoldberg2012}. \citet{BhansaliEtAl2024} show that treating trend following as a
separate asset class within risk parity improves its performance, notably in 2022, when stocks and bonds
fell together. Return-stacking funds bring levered diversification to retail investors by combining a
full exposure to stocks or bonds with futures-based diversifiers \citep{GordilloHoffstein2021}.

Our diversifying asset classes differ in their economic sources of return. Trend-following strategies
earn positive average returns across asset classes and centuries and have historically performed well
in extended equity declines \citep{MoskowitzOoiPedersen2012,HurstOoiPedersen2017}; their returns are
commonly attributed to initial underreaction and delayed overreaction to news, not to leverage aversion.
Gold is a store of value whose real price depends on the monetary regime \citep{ErbHarvey2013}. Studies
of retirement withdrawals find that trend following can reduce sequence-of-returns risk
\citep{ClareSeatonSmithThomas2017}.

\section{Data}
\label{sec:data}

Following \citet{AnarkulovaCederburgODoherty2022,AnarkulovaCederburgODoherty2025}, we model
forward-looking returns with the historical returns of a broad cross section of developed countries.
All inputs to the simulations come from public sources.\footnote{Proprietary indexes of commercial
trend-following funds are used only to validate the managed-futures series.} The starting point is release 6 of the Jord\`a-Schularick-Taylor
Macrohistory Database \citep{JordaEtAl2019}, which contains annual nominal total returns on stocks and
ten-year government bonds, short-term bill rates, consumer price indexes, and exchange rates for 16
developed countries through 2020. We extend each series through 2025 using dividend-adjusted local
indexes, locally listed exchange-traded funds, and the Global Macro Database
\citep{MullerXuLehbibChen2025}. The dataset is an unbalanced panel of 1,561 country-years from 1927 to
2025. Section~\ref{app:data-reconstruction} of the internet appendix provides details on the data sources
and return calculations.

All returns for a given country-year are real returns for local investors. The international stock
return for a country is the market-capitalization-weighted return on all other markets, converted into
local currency, with weights from \citet{KuvshinovZimmermann2022}. As in \citet{AnarkulovaCederburgODoherty2025}, the stock allocation is 33\% domestic stocks
and 67\% international stocks, and we refer to this mix as stocks.

Global bonds are an equally weighted portfolio of the ten-year government bonds of all countries in the
panel in a given year, including the investor's own country. Foreign bond positions are hedged into local currency
with one-period forwards on the initial principal, priced at covered interest parity. We deduct a fund fee and a hedging cost of 0.10\% per year each.\footnote{Liquid forward markets did not exist for many panel currencies
before the 1970s, and capital controls restricted cross-border positions under the Bretton Woods system.
Hedged returns before then are therefore a hypothetical construction. The covered-parity assumption also
ignores deviations such as the cross-currency basis documented by \citet{DuTepperVerdelhan2018}.}

Gold returns are based on the annual US dollar price of gold from MeasuringWorth
\citep{OfficerWilliamsonMeasuringWorth} through 1967 and the London Bullion Market Association afternoon
fixing \citep{LBMA} from 1968, when the private London market was separated from the official price,
less a custody cost of 0.40\% per year. Before 1968, the dollar price of gold
was administered, so early gold returns reflect a different monetary regime \citep{ErbHarvey2013}, and we
examine the post-1970 period separately.

No investable managed-futures index spans our sample period, so we construct a time-series momentum
strategy from monthly futures-equivalent returns following \citet{MoskowitzOoiPedersen2012} and
\citet{HurstOoiPedersen2017}. The signal for each market averages the signs of its prior one-, six-, and
twelve-month returns, lagged one month. Markets are scaled to 15\% volatility and sectors weighted by
inverse volatility; the portfolio is scaled to 10\% volatility over the preceding 36 months, with leverage
capped at three. The universe expands as data become available across
equity indexes, government bonds, commodities, and currencies. The strategy earns US dollar collateral
income, and we deduct a management fee of 0.85\% per year and trading costs of three basis points per
unit traded, or about 0.61\% per year. The annual series is hedged into local currency in the same manner
as global bonds.

The managed-futures series has important limitations. It omits commodity carry and roll returns, uses
synthetic bond futures and price-only equity indexes, and trades a narrow set of markets before 1960;
currencies enter only in 1971. Its signal follows studies published after much of our sample, so its
design is not free of hindsight. For a US investor, the series earns an average real return of 8.05\% per
year with a standard deviation of 14.18\% over 1927 to 2025 (Table~\ref{tab:mf-diagnostics} of the
internet appendix). The most direct test of the series is its record against commercial funds. Over 2000
to 2025, it earned 6.11\% per year in US dollars, compared with 4.20\% for the SG CTA index, 5.38\% for the
SG Trend index, and 4.09\% for the BTOP50 index, and its Sharpe ratio was 0.41, compared with 0.27 to 0.30.
The gap comes from 2000 to 2009. Over 2010 to 2025, the series earned 2.44\% per year with a Sharpe ratio
of 0.14, below all three indexes (2.70\% to 3.64\% and 0.18 to 0.24), consistent with the weaker
trend-following returns documented by \citet{BabuEtAl2020}. The series is also only moderately
correlated with the indexes, at 0.55 to 0.59 over 2000 to 2025. Section~\ref{sec:robustness} reports the
reduction in managed-futures returns at which our conclusions change.

Table~\ref{tab:sleeves} summarizes the properties of the four asset classes. Over the full sample, stocks
earn the highest average real return (8.12\% per year), followed by managed futures (7.37\%), gold
(3.99\%), and global bonds (1.94\%). In excess of the local bill return, managed futures have a Sharpe
ratio of 0.49, compared with 0.31 for stocks, 0.29 for bonds, and 0.12 for gold. The rankings are the same
after 1970. In the country-years in the worst decile of stock returns, stocks lose 25.0\% on average,
bonds lose 2.1\%, gold gains 2.6\%, and managed futures gain 8.4\%. After 1970, gold and managed futures
gain 10.5\% and 15.4\% in these years. Bonds suffer during inflationary periods, with a ten-year
correlation with inflation of $-0.78$, similar to the $-0.78$ reported by
\citet{AnarkulovaCederburgODoherty2025} for 30-year horizons.\footnote{The annual correlation between gold
and stocks in the full sample is 0.60. This correlation reflects country-years with large currency
depreciations and high inflation, in which both gold and international stocks rise in local currency. At
the ten-year horizon, the correlation is $-0.35$.}
\exhibit{figures/erc/sleeves.tex}

\section{Methods}
\label{sec:methods}

We adopt the lifecycle framework of \citet{AnarkulovaCederburgODoherty2025} at an annual frequency.
Their analysis shows that a fixed-weight strategy of 33\% domestic and 67\% international stocks achieves
virtually the same expected utility as their optimal age-based strategy, so we focus on fixed-weight
strategies throughout.

\subsection{Lifecycle design and income}
\label{sec:design}\label{sec:model}\label{sec:income}\label{sec:methods-household}

Households are composed of a female and a male of equal age. Each member of the household is eligible to
work and save from age 25 to age 64, and the couple retires at age 65. In the base case, individuals save
$r_c=10\%$ of their labor income for retirement, and individuals earning less than $Y_{\min}=\$15{,}000$
(in 2022 USD) in a given year do not contribute. Contributions are made at the beginning of each year. We
model labor income using model 6 of \citet{GuvenenKarahanOzkanSong2021}, as in
\citet{AnarkulovaCederburgODoherty2025}. Log income contains a quadratic age profile, a persistent
component with autoregressive coefficient 0.959, transitory shocks, and age- and income-dependent
nonemployment. Both members of the couple have the median parameter values. Social Security benefits are
computed from each simulated earnings history using the 2022 benefit formulas, including spousal and
survivor benefits, and Supplemental Security Income provides a consumption floor in retirement.

At retirement, the couple withdraws $r_w=4\%$ of its retirement wealth in the first year and the same
real amount in each subsequent year \citep{Bengen1994}. The fixed real withdrawal makes retirement
outcomes sensitive to the order of returns \citep{Estrada2018}. Mortality follows sex-specific
Gompertz-Makeham laws calibrated to the longevity moments reported by
\citet{AnarkulovaCederburgODoherty2025} from Social Security Administration tables \citep{SSA2020}, and
the lifespan of each individual is random.\footnote{As in \citet{AnarkulovaCederburgODoherty2025}, the
model combines US earnings, public pensions, and mortality with returns pooled across developed
countries. The simulations therefore describe a US household exposed to the return histories of other
countries. They do not model the pension systems of those countries.}

\subsection{Household utility and allocation strategies}
\label{sec:utility}\label{sec:portfolio-rules}

Household utility is determined by annual retirement consumption and a bequest:
\begin{equation}
 U(C,B)=\sum_{t=65}^{T_{\max}}
 \frac{(C_t/\sqrt{H_t})^{1-\gamma}}{1-\gamma}
 +\theta\frac{(B+k)^{1-\gamma}}{1-\gamma},
 \label{eq:utility}
\end{equation}
where $C_t$ is real consumption at age $t$, $H_t$ is the household's size, $B$ is the real bequest, and
$T_{\max}$ is the age at death of the last survivor. Following
\citet{AnarkulovaCederburgODoherty2025}, we use the estimates of \citet{DeNardiFrenchJones2010}: a
coefficient of relative risk aversion of $\gamma=3.84$, a bequest intensity of $\theta=2{,}360$, and a
bequest curvature parameter of $k=\$490{,}000$. \citet{AnarkulovaCederburgODoherty2025} evaluate monthly
consumption and multiply the annual bequest intensity by $12^\gamma$. With uniform consumption within
the year, our annual specification preserves the relative weight of consumption and bequest in their
utility function.

The household invests in stocks, global bonds, gold, and managed futures and can borrow at the local bill
rate plus a spread. Let $w$ be the vector of asset class weights as fractions of wealth, and let
$G=\mathbf{1}'w$ be gross exposure. The real portfolio return in year $t$ is
\begin{equation}
 R^p_t=w'R_t+(1-G)R^f_t-\max(G-1,0)\,\phi-h\kappa,
 \label{eq:return}
\end{equation}
where $R_t$ is the vector of real asset class returns, $R^f_t$ is the local real bill return, $\phi$ is
the borrowing spread, $h$ is the weight of bonds and managed futures, and $\kappa$ is the hedging cost. In the base
case, $\phi=0.30\%$, which is close to the lowest spread of 0.37\% considered by
\citet{AnarkulovaCederburgODoherty2025} and estimated by \citet{vanBinsbergenDiamondGrotteria2022} from
derivative prices. This is the cost of financing through box spreads on index options \citep{WangBoxFinancing2025} or
through the futures held by return-stacked funds, both open to individual investors
(Section~\ref{sec:implementation}). It is not the rate of a retail margin loan, which we cover with higher
spreads.

We consider three families of strategies, each evaluated at gross exposures of $G=100\%$, 125\%, 150\%,
175\%, and 200\%:
\begin{itemize}[leftmargin=2em,itemsep=2pt]
\item the \emph{all-equity strategy}, which invests $G$ in stocks;
\item the \emph{proportional strategy}, which invests $0.40G$ in stocks, $(4/15)G$ in global bonds, $G/6$
in gold, and $G/6$ in managed futures; and
\item the \emph{risk parity strategy}, which invests $0.1453G$ in stocks, $0.4747G$ in global bonds,
$0.1205G$ in gold, and $0.2595G$ in managed futures.
\end{itemize}
The proportional weights start from a portfolio of 60\% stocks and 40\% bonds and borrow 50\% of wealth
to add 25\% gold and 25\% managed futures; the resulting proportions are then held fixed as gross exposure
varies. This practitioner rule is similar to the return-stacking portfolios offered to retail investors,
and its weights were fixed before we examined household outcomes. We keep it as our main diversified
strategy even though bonds turn out to lower utility, because dropping an asset class after observing its
contribution would select on outcomes. We also report a strategy that holds only stocks and managed
futures in the same relative proportions, and the weights that maximize expected utility in our sample.
The risk parity weights equalize the contributions of the four asset classes to the variance of an
unlevered portfolio \citep{MaillardRoncalliTeiletche2010}, using the sample covariance matrix of annual
real returns over all 1,561 country-years. This estimation uses the full sample and therefore gives the
risk parity strategy an informational advantage.

Gross exposure is not a measure of risk: at 175\%, the three families have annual volatilities of 44.0\%,
26.6\%, and 17.1\%. Our main comparison therefore scales each diversified strategy to the 25.28\% annual
volatility of the unlevered all-equity strategy in the pooled panel, the comparison implied by the
leverage-aversion argument \citep{Asness2024}. This gives exposures of 166\% for the proportional strategy,
265\% for risk parity, and 132\% for stocks and managed futures, which we hold fixed in every sensitivity
analysis.

Gross exposure $G$ is the sum of the allocations to the four asset classes as fractions of wealth, and
$G-1$ is borrowing. It is not the notional value of the underlying contracts. The managed-futures series is
itself a futures strategy that targets 10\% volatility and can hold up to three times its capital in
notional positions, and bonds and foreign returns are hedged with forwards. A 28\% allocation to managed
futures counts as 28\% of wealth whatever its internal notional. The borrowed part of the portfolio
requires collateral and can be liquidated after losses; Section~\ref{sec:robustness} reports margin tests
and Section~\ref{sec:implementation} discusses financing.

\subsection{Simulation procedure}
\label{sec:simulation}\label{sec:bootstrap}

We simulate 10,000 lifecycle outcomes. Each draw combines random lifespans for the two members of the
couple, simulated labor incomes, and a time series of annual real asset class returns from a stationary
block bootstrap in the spirit of \citet{PolitisRomano1994}. We select a random starting country-year and
draw a block of consecutive years from the same country, with a geometric block length of ten years on
average; each country-year supplies all asset class returns jointly. If a country's data end before the
block is complete, the block continues at the start of a randomly drawn country's sample, without a
currency conversion. Wealth evolves as $W_{t+1}=\max\{0,(W_t+S_t)(1+R^p_{t+1})\}$ from $W_{25}=0$, where
$S_t$ is the contribution and $R^p$ is given by equation~\eqref{eq:return}. In retirement, the couple
withdraws $D_t=\min\{0.04\,W_{65},W_t\}$ at the beginning of each year and consumes
$C_t=\max\{D_t+SS_t,SSI_t\}$, where $SS_t$ is the Social Security benefit and $SSI_t$ is Supplemental
Security Income. The bequest is the wealth left at the death of the last survivor. Expected utility is the
average of equation~\eqref{eq:utility} across draws. All strategies share the same return, income, and
mortality draws, and we find equivalent savings rates by bisection. Portfolios are rebalanced to their
target weights once a year.

\subsection{Three measures of retirement outcomes}
\label{sec:outcomes}

We compare strategies with three measures, summarized in Table~\ref{tab:measures}. The first is the
equivalent savings rate of \citet{AnarkulovaCederburgODoherty2025}, the savings rate $r_c^*$ that gives a
strategy the same expected utility as the all-equity strategy with 100\% exposure and a 10\% savings rate:
\begin{equation}
 E[U_j\mid r_c^*]=E[U_{\text{All-equity}}\mid r_c=10\%].
 \label{eq:eqsaving}
\end{equation}
A rate below 10\% means that the strategy provides higher expected utility. It measures utility, not
income: a household saving $r_c^*$ has the same expected utility as the all-equity household, not the same
retirement income on each path.

The second is the probability of financial ruin, defined as exhausting retirement wealth under the 4\% rule
before the death of the last survivor. Contributions, retirement wealth, and the withdrawal all scale with
the savings rate, so this probability does not depend on the savings rate
(Section~\ref{app:framework} of the internet appendix). A strategy with a low expected return can
therefore have a low ruin probability because it withdraws a small income.

The third measures the ability to fund a given income. Each household withdraws 4\% of the retirement
wealth that the all-equity strategy with 100\% exposure and a 10\% savings rate accumulates on the same
simulated path, whatever strategy it follows. We call this withdrawal the path-matched income: it differs
across paths and is high when equities have done well. The path-matched income failure is the probability
that the strategy's wealth cannot pay it before the death of the last survivor. We compute it with a 10\%
savings rate and, in Table~\ref{tab:measures}, also at the strategy's own equivalent savings rate.
\exhibit{figures/erc/measures.tex}

\section{Results}
\label{sec:results}\label{sec:leverage-ladders-main}

\subsection{Levered diversification}
\label{sec:main}

Table~\ref{tab:central-results} reports the asset class weights, return volatility, ruin probability, and
equivalent savings rate for each strategy. Our annual implementation closely matches
\citet{AnarkulovaCederburgODoherty2025}: the unlevered all-equity strategy has a ruin probability of
7.0\%, as in their study, and their ruin probabilities of 17.1\% for domestic stocks and 16.9\% for a 60/40
strategy compare with 15.5\% and 13.7\% in our data, although their data cover 39 countries from 1890.
\exhibit{figures/erc/central.tex}

Diversification and leverage improve expected utility only in combination. Panel B shows that the
unlevered proportional strategy lowers the ruin probability from 7.0\% to 5.0\% but raises the equivalent
savings rate to 12.35\%: it reduces risk but gives up too much expected return. Panel A shows that leverage
without diversification provides small utility gains at a large cost in capital preservation. The
equivalent savings rate of the all-equity strategy falls to 8.56\% at 125\% exposure and 8.30\% at 150\%,
but rises to 9.64\% at 175\% and 15.66\% at 200\%, while its ruin probability rises at every step, from 7.0\%
to 20.3\%.

Panel E reports the strategies at equal volatility. The proportional strategy at 166\% exposure has an
equivalent savings rate of 6.35\% and a ruin probability of 2.3\%; the 95\% confidence interval for its
difference in ruin probability with the all-equity strategy is $-5.2$ to $-4.3$ percentage points. Its
returns are as volatile as those of the all-equity strategy but much less concentrated in stock market
outcomes. On the exposure grid, the gains increase with leverage up to 200\%, where the equivalent savings
rate is 4.72\% and the ruin probability 1.8\%, at a volatility of 30.3\%, above that of the all-equity
strategy. The stocks-and-managed-futures strategy reaches equity volatility at only 132\% exposure, with an
equivalent savings rate of 5.22\% and a ruin probability of 2.1\%.

Risk parity illustrates why the basis of comparison matters. At 175\%, where its volatility is only
17.1\%, it does worse than the proportional strategy (8.29\% equivalent savings rate, 3.8\% ruin). At
equal volatility, it needs 265\% exposure and does better (4.18\% and 1.9\%). It attains this result by borrowing 165\% of wealth and by
holding 69\% of wealth in managed futures, and Section~\ref{sec:robustness} shows that its advantage is
fragile.

\subsection{Capital preservation and income funding}
\label{sec:common-income}

Table~\ref{tab:common} and Figure~\ref{fig:income-ruin} compare the probability of financial ruin with the
path-matched income failure on the exposure grid, with a 10\% savings rate. The two outcomes diverge
sharply for strategies with low expected returns. The unlevered proportional strategy has a lower ruin
probability than the all-equity strategy (5.0\% versus 7.0\%), but it fails to fund the path-matched
income in 31.5\% of lifetimes. The unlevered risk parity strategy fails in 47.4\%. These strategies
preserve capital by accumulating less of it. The levered all-equity strategy shows the opposite pattern:
at 175\% exposure, its ruin probability is 14.8\%, but it fails to fund the path-matched income in only
6.3\% of lifetimes.
\exhibit{figures/erc/common.tex}
\exhibit{figures/erc/fig_income_ruin.tex}

At 175\% exposure, the risk parity strategy has a ruin probability of 3.8\% but a 14.2\% probability of
failing to fund the path-matched income. With nearly half of its unlevered weight in bonds, it produces
stable but small retirement accounts. At equal volatility, the three diversified strategies fund the
path-matched income more reliably than the all-equity strategy when they save 10\%, with failure
probabilities of 3.2\% for the proportional strategy, 2.2\% for risk parity, and 0.8\% for stocks and
managed futures (Table~\ref{tab:measures}).

The comparison changes when each strategy saves only its equivalent savings rate. The proportional
strategy then fails to fund the path-matched income in 13.1\% of lifetimes and risk parity in 16.9\%,
compared with 7.0\% for the all-equity strategy. About 70\% of these failures occur in paths where
the all-equity strategy accumulates more than its median wealth, so that the income to fund is high: the
diversified strategies give up part of the upside of equity booms in exchange for protection in equity
crises, which is what the utility function rewards. Only the stocks-and-managed-futures strategy funds the
path-matched income more reliably than the all-equity strategy at its own equivalent savings rate (5.9\%
failure). Capital preservation, utility, and income funding are therefore distinct objectives, and a low
ruin probability can reflect a low planned income rather than a safe portfolio.

\subsection{Sources of the gains}
\label{sec:sleeves}

The gains come mainly from managed futures, and bonds lower utility when they replace stocks.
Table~\ref{tab:composition-value} reports changes in the equivalent savings rate at 175\% exposure, the
grid point closest to equal volatility for the proportional strategy, in the full sample and after 1970.
Panel A adds each diversifying asset class to the levered all-equity strategy at its proportional-strategy
weight, financed by reducing stocks. Panels B and C remove each asset class from the diversified
strategies and reallocate its weight to stocks. Table~\ref{tab:ablations} of the internet appendix
reports ruin and path-matched income outcomes in the full sample.
\exhibit{figures/erc/composition_value.tex}

Managed futures are indispensable. Adding them to the levered all-equity strategy lowers the equivalent
savings rate by 5.4 percentage points in the full sample and 8.3 points after 1970. Removing them raises the
equivalent savings rate of both diversified strategies in both samples, and, with their weight spread pro
rata (Table~\ref{tab:ablations}), ruin from 2.1\% to 3.7\% (proportional) and 3.8\% to 8.1\% (risk parity).

In our panel and at a given exposure, global bonds lower utility. Removing them lowers the equivalent
savings rate of the proportional strategy by 1.7 and 1.0 points in the two samples and, spread pro rata,
its path-matched income failure from 2.2\% to 0.9\%. Panel D of Table~\ref{tab:central-results} reports two
bond-free versions of the proportional strategy at 175\%: with stocks, gold, and managed futures, the
equivalent savings rate is 4.15\%, and with stocks and managed futures only, 3.25\%. These results are
consistent with the long-horizon properties of bonds documented by \citet{AnarkulovaCederburgODoherty2025}.
They concern hedged ten-year government bonds in a panel that includes the inflations of the 1940s and
1970s, and bonds can still add value indirectly, by lowering volatility and so permitting more leverage on
the other asset classes, as in the risk parity strategy at equal volatility.

Gold's contribution depends on the monetary regime. In the full sample, which includes four decades of
administered gold prices, removing gold from the proportional strategy lowers its equivalent savings rate
by 0.8 points. After 1970, adding gold to the levered all-equity strategy lowers the equivalent savings
rate by 7.3 points, nearly as much as managed futures, and removing it from the proportional strategy
raises the rate by 0.4 points. Once managed futures are present, gold adds little, as the two asset
classes pay off in the same episodes.

Panel D of Table~\ref{tab:composition-value} reports the weights that maximize expected utility at 175\%
exposure, an in-sample upper bound on the gains. In the full sample, the optimal strategy lowers the
equivalent savings rate by 6.9 points relative to the levered all-equity strategy, compared with 3.8
points for the proportional strategy. The optimal strategy allocates 92\% of wealth to managed futures and
none to bonds or gold. When managed-futures returns are reduced by 3\% or 6\% per year, the optimal weights
shift toward stocks and gold, and the gain shrinks to 4.6 and 3.5 points. The proportional rule captures
about half of the in-sample gain, and removing bonds or gold would improve it, which suggests that it was
not tuned to our sample.\footnote{Section~\ref{app:local-decomposition} of the internet appendix decomposes
these changes into a mean component and a covariance component with the marginal utility of return. For
small changes in weights, the covariance component accounts for most of the value of adding each asset
class to the levered all-equity strategy.}

\subsection{Historical uncertainty and sample periods}
\label{sec:history}

The 1,561 country-years of the panel are not 1,561 independent observations. Gold and managed futures are
global series, common to all countries in a year up to currency conversion, and stock and bond returns
are correlated across countries. For these asset classes, the independent information is closer to the
99 calendar years of the sample. The confidence intervals of Section~\ref{sec:main} measure only Monte Carlo
error given the panel. To measure uncertainty about the history itself, we resample blocks of complete
calendar-year cross-sections, with average lengths of 5, 10, or 20 years, and repeat the simulation on
each of 100 resulting histories. Table~\ref{tab:restored-history} reports the results at equal volatility.
The proportional strategy has a lower ruin probability than the all-equity strategy in 97\% to 100\% of the
histories, a lower equivalent savings rate in 96\% to 99\%, and a lower path-matched income failure in
78\% to 83\%. The stocks-and-managed-futures strategy does better on all three measures in 99\% to 100\% of
the histories.
\exhibit{figures/erc/history.tex}

Panel B of Table~\ref{tab:sensitivity} restricts the sample to 1970 to 2025. The ruin probabilities of the
diversified strategies fall to 0.5\% or less, compared with 6.7\% for the all-equity strategy. These results
require caution. The sample begins just before 1973 and 1974, the strongest joint episode for gold and
trend following in our data (Table~\ref{tab:mf-diagnostics}), and excludes the wartime inflations and
currency collapses of the 1930s and 1940s. We therefore interpret the post-1970 results as evidence for
the period of free gold markets and broader managed-futures coverage, not as an out-of-sample test.

Panel B also excludes Italy in 1942, when the lira moves from 19.7 to 225.5 per US dollar in the
Macrohistory Database, giving an Italian investor a real return of about 1,000\% on world stocks, by far
the largest observation in the panel. The results are essentially unchanged. Restricting the proportional
strategy to asset classes that were investable at each date, with no gold before 1975 and
managed-futures sectors added as their futures markets opened, gives a ruin probability of 2.3\% and an
equivalent savings rate of 5.77\% at 175\% exposure (Section~\ref{app:availability} of the internet
appendix).

\subsection{Borrowing costs, managed-futures returns, and margin requirements}
\label{sec:robustness}\label{sec:costs}\label{sec:mf-proxy}

Table~\ref{tab:sensitivity} reports the three measures at equal volatility under alternative
assumptions.
\exhibit{figures/erc/sensitivity.tex}

\subsubsection{Borrowing costs}

Panel C raises the borrowing spread from 0.30\% to 1.00\%, 2.00\%, and 3.00\%. The equivalent savings
rate of the proportional strategy rises to 7.21\%, 8.66\%, and 10.45\%. By linear interpolation, it loses
its saving advantage over the unlevered all-equity strategy at a spread of about 2.75\%, and its advantage
in funding the path-matched income at about 1.3\%. At the medium spread of 1.40\% of
\citet{AnarkulovaCederburgODoherty2025}, which equals the lowest broker margin rate in April 2024, its
interpolated equivalent savings rate is 7.79\%. Households need not pay margin rates: box spreads on
S\&P 500 options have traded at a median of 0.20\% to 0.31\% above Term SOFR
\citep{CMEBoxFinancing2024}, and the S\&P 500 financing spread embedded in futures averaged 0.3\% over
2021 to 2023, although it peaked above 1.4\% in December 2024 \citep{DEShaw2025}. Risk parity, which borrows most, loses its saving advantage
at a spread of about 2.1\%. The stocks-and-managed-futures strategy, which borrows only 32\% of wealth,
keeps all three advantages at a 3.00\% spread, with an equivalent savings rate of 6.62\% and a
path-matched income failure of 1.7\%.

\subsubsection{Managed-futures returns}

Panel D subtracts an additional 2\%, 3\%, or 6\% per year from managed-futures returns, beyond fees and
trading costs. With a 2\% reduction, the proportional strategy has an equivalent savings rate of 7.39\%, a
ruin probability of 2.8\%, and a path-matched income failure of 6.0\%, all better than the all-equity
strategy. It loses its saving advantage at a reduction of about 5.9\% per year and its advantage in funding
the path-matched income at about 2.5\%. The stocks-and-managed-futures strategy keeps
its saving advantage up to about 5.9\% and its income advantage up to about 4.7\%. Risk parity is the most
sensitive, losing its income advantage at about 1.5\%.

Panel E rebuilds the managed-futures series with other signals: a 1/3/12-month blend, the 12-month
signal alone, the 1-month signal alone, and the baseline blend with a fee of 1.70\% instead of 0.85\%. The
alternative series have annual correlations of 0.76 to 0.96 with the baseline. The proportional strategy's
equivalent savings rate ranges from 6.29\% to 7.16\%, its ruin probability from 2.2\% to 2.4\%, and its
path-matched income failure from 2.7\% to 5.1\%, all better than the all-equity strategy. The results do not depend on the signal.

\subsubsection{Margin requirements and other design choices}

Our model is annual, and it does not simulate margin calls within the year. In the simulations, the
proportional strategy at 200\% exposure never breaches a 25\% maintenance requirement at year-end, whereas
the all-equity strategy breaches it in 2.1\% of account-years at 200\% (Table~\ref{tab:restored-margin}).
In monthly US data from 1970 to 2025, the proportional strategy never breaches a 25\% or 40\% requirement
at month-end at any exposure up to 200\%, while the all-equity strategy breaches the 40\% requirement in 13
years at 200\% (Table~\ref{tab:monthly-margin}). We have not tested the risk parity strategy at 265\%
exposure. The internet appendix also shows that, at 200\% exposure, the advantage of the levered
proportional strategy holds for risk aversion from $\gamma=2$ to $\gamma=10$ (Table~\ref{tab:restored-gamma}), savings rates from 5\% to
15\% and withdrawal rates from 3\% to 5\% (Table~\ref{tab:restored-policy}), and, at 175\%, in every quintile
of the ex post real bill rate (Table~\ref{tab:bill-quintiles}).\footnote{We take a partial equilibrium view. Large
flows into gold would require price adjustments \citep{ErbHarvey2013}, inelastic demand can make equity
prices sensitive to flows \citep{GabaixKoijen2023}, and crowded trend-following trades can raise execution
costs \citep{BaltasCrowding2019}.}

\section{Implementation}
\label{sec:implementation}

Individual investors can hold these strategies today; we verify it for a European household, using
instruments available in September 2026. UCITS funds may reach 200\% gross exposure
\citep{ESMAGlobalExposure2010}. The WisdomTree Global Efficient Core UCITS ETF holds 90\% global equities
and 60\% government bond futures for 0.25\% a year \citep{WisdomTreeNTSG2026}, and the iMGP DBi Managed
Futures UCITS ETF, listed in Paris since 2025, charges 0.75\% \citep{iMGPDBMF2025}. With physically backed
gold, putting 74\% of wealth in the first fund and 13\% in each of the other two gives 66\% stocks, 44\%
bonds, 13\% gold, and 13\% managed futures, a gross exposure of 137\% financed inside the funds at
futures-implied rates. The rest of the leverage, 29\% of wealth for the proportional strategy at equal
volatility and 32\% for stocks and managed futures, can be borrowed through box spreads on index options,
available in ordinary brokerage accounts and priced at 0.20\% to 0.50\% above SOFR or Treasury yields
\citep{CMEBoxFinancing2024,WangBoxFinancing2025}. Our returns are before personal taxes.

\section{Conclusion}

We study the lifecycle investment problem of \citet{AnarkulovaCederburgODoherty2025} adding global bonds, gold, managed futures, and borrowing to the opportunity set. Diversification alone lowers the probability of ruin but not expected utility, and leverage alone adds
little utility and raises ruin. Together, they improve all three measures. Levered to the
volatility of unlevered equities, a multi-asset strategy matches the all-equity utility with a 6.35\%
savings rate instead of 10\%, lowers the probability of ruin from 7.0\% to 2.3\%, and, at the same savings
rate, funds the all-equity income more reliably. The gains come mainly from managed futures, which pay off in equity crises, rather than from levering
bonds, and stocks with managed futures alone do better still. Equal utility is not equal income, preserving capital is not funding an income, and
rankings differ at equal exposure and equal risk. With financing close to the bill rate, levering assets that pay off in equity crises is safer than concentrating in stocks.

\ifsubmission\else\clearpage\fi
\singlespacing
\printbibliography[title={References}]

\ifsubmission
\clearpage
\exhibitlist
\fi

\end{document}
